\documentclass[11pt]{article}
\usepackage{amsmath,amssymb,amsthm,hyperref}
\newtheorem{proposition}{Proposition}
\title{Lookup Identities and Hyperplane Cover Stability}
\author{}
\date{November 23, 2026}
\begin{document}
\maketitle

\paragraph{Scope.}
This note audits three additional literature mechanisms. None supplies
the universal head bound. The pole obstruction below concerns equality
of rational functions on an interval. It is not a lower bound for strict
sign representation on a finite Boolean cube.

\section{Logarithmic-derivative lookup identities}

Hab\"ock's \emph{Multivariate lookups based on logarithmic derivatives},
\href{https://eprint.iacr.org/2022/1530}{Cryptology ePrint 2022/1530},
Section 3, equations (13)--(15), uses the following identity. For two
finite lists of field elements $a_i$ and $t_j$, with distinct table
entries $t_j$, membership of every $a_i$ in the table is certified by
\[
 \sum_i\frac{1}{X+a_i}=\sum_j\frac{m_j}{X+t_j},
 \qquad m_j=\#\{i:a_i=t_j\}.
\]
In positive characteristic the characteristic must be large enough to
preserve the relevant multiplicities. The protocol supplies multilinear
table, witness, and helper functions and tests a polynomial identity.
These functions are not required to be affine.

The source's use of a Boolean hypercube indexes a table of field
elements. It does not make that table an affine function of its index.
Likewise, testing the identity at a verifier's random challenge is not
the same operation as producing the label of an arbitrary cube input.

\begin{proposition}[Distinct lookup poles cannot disappear in a flat identity]
Let $t_1,\ldots,t_M$ be distinct real numbers and let each $m_j$ be
nonzero. Suppose that on a nonempty open real interval, avoiding all
displayed denominator zeros,
\[
 \sum_{j=1}^M\frac{m_j}{X+t_j}
 =c+\sum_{h=1}^H\frac{a_hX+b_h}{u_hX+v_h}.
\]
Assume that no denominator on the right is the zero polynomial.
Then $H\ge M$.
\end{proposition}
\begin{proof}
Multiplication by all denominators gives a polynomial identity on an
open interval, hence a polynomial identity everywhere. Thus the two
rational functions agree. Fix $j$. The left side has a simple pole
at $-t_j$, because its residue is $m_j\ne0$ and all other left-hand
summands are regular there. If every right-hand denominator were
nonzero at $-t_j$, the right side would be regular there, a
contradiction. Consequently some right-hand denominator vanishes at
$-t_j$. A nonzero affine polynomial vanishes at at most one real
number. Distinct $j$ therefore require distinct right-hand denominators,
and $H\ge M$.
\end{proof}

The proposition permits cancellation between right-hand summands,
constant denominators, and arbitrary affine numerators. Cancellation
can remove a pole, but cannot create a pole at a new location.
It does not prevent a different, target-dependent score from matching
only the signs at finitely many points. No such compressed score is
constructed by the lookup identity alone.

\section{Functional lower bounds for reciprocal polynomials}

Hakoniemi, Limaye, and Tzameret,
\emph{Functional Lower Bounds in Algebraic Proofs: Symmetry, Lifting,
and Barriers},
\href{https://www.doc.ic.ac.uk/~itzamere/HLT-functional-23.pdf}{March 2024 manuscript},
Section 6.2 and Claim 55, use alternating sums of restrictions of an
inverse knapsack polynomial to analyze its coefficient spaces on
Boolean inputs. In their notation the relevant identity is
\[
 h_\alpha=\sum_{\mu\subseteq\alpha}(-1)^{|\mu|}
 \tau_\mu(1/\mathrm{ks}_w).
\]
Their word-indexed knapsack polynomial $\mathrm{ks}_w$ includes products
of input variables. An inverse of this polynomial is not generally an
allowed head. Even the displayed alternating expansion retains one
term for each subset of $\alpha$. The paper proves lower bounds for
polynomial proof representations; it does not give a head-efficient
compiler for these inverses or their sign tables.

\section{The new cover-stability theorem}

Das, Djaljapayan, Lin, and Yu,
\emph{Stability for Hyperplane Covers},
\href{https://www.combinatorics.org/ojs/index.php/eljc/article/download/v33i4p37/pdf}{Electronic Journal of Combinatorics 33(4), P4.37},
published November 20, 2026, sharpen the multiplicity-cover problem
that avoids the origin. Define
\[
 w(x)=\frac{1}{|x|\binom n{|x|}},\qquad
 w(h)=\sum_{x\in h\cap\{0,1\}^n}w(x),
\]
where $h$ avoids the origin. Their Theorem 3 characterizes weight-one
hyperplanes $\sum_i c_i x_i=1$: the coefficients are integers at most
one and their sum is at least one. Otherwise $w(h)\le1-1/n$.
Corollary 4 counts $\binom{2n-1}{n}$ weight-one hyperplanes.

These are weighted incidence and stability statements for covers of
all nonzero vertices with prescribed multiplicity. They do not cover
an arbitrary target subset while avoiding its entire complement, and
they do not construct fractional scores with arbitrary labels. In
particular, the stability theorem does not remove the factor-$n$ gap
in the general exact-cover bounds used in the earlier audit. A
conversion of cover unions into strict-sign sums would also still
need a separate amplitude argument.

\paragraph{Audit outcome.}
No denominator list, numerator list, or parameter construction at
$O(2^n/n^2)$ heads follows from the three mechanisms examined here.
The exact manuscript model and the universal target remain unchanged.

\end{document}
